\section{Discussion}
The results confirm that tree-structured memory provides clear advantages for questions requiring multi-step or temporal reasoning, while flat methods retain an edge on open-domain recall. The 8.6-point F1 gain on multi-hop and 14.1-point gain on temporal questions over Memory~R1 suggest that the hierarchical grouping of facts by speaker, topic, and session enables retrieval paths that flat scoring cannot replicate. At the same time, the dominance of RAG~(BM25) on open-domain questions (39.1 F1) highlights that broad keyword coverage remains valuable when questions do not require structural navigation.

The consistent F1--J divergence across baselines is also informative. Methods that use structured memory (Semantic XPath, our method) tend to achieve higher J relative to their F1, suggesting that tree-guided retrieval surfaces more contextually complete evidence even when the final answer does not lexically match the gold reference. This pattern reinforces the value of evaluating with both lexical and semantic metrics.

The current formulation uses oracle structured data for memory construction, meaning the Memory Manager has not yet been trained end-to-end. This represents an important caveat: the reported gains reflect the ceiling of what tree structure can provide given perfect writing, not what an end-to-end system would achieve. Completing the Memory Manager training loop is the most important next step.

More broadly, the results suggest that structural ambition must be balanced against training stability. The present method uses stable CRUD updates for memory writing, while the Retrieve Agent carries the burden of structure-aware navigation. This division is cleaner than asking a single policy to both reorganize memory arbitrarily and discover retrieval strategies simultaneously.
